\subsection{Copy constraints in PLONK}
\label{copy-constraints}

Now we elaborate on Step 3 which we deferred back in
{[}copy-constraint-deferred{]}. As an example, the constraints might be:
\[a_{1} = a_{4} = c_{4}\quad\text{and }\quad b_{2} = c_{1.}\] Before we
show how to check this, we provide a solution to a ``simpler'' problem
called ``permutation check''. Then we explain how to deal with the full
copy check.

\subsubsection{Easier case: permutation check}

strong(Problem):\\

Well, actually, it would be necessary and sufficient for the identity
\[\begin{array}{r}
\left( T + P\left( \omega^{1} \right) \right)\left( T + P\left( \omega^{2} \right) \right)\ldots\left( T + P\left( \omega^{n} \right) \right) \\
 = \left( T + Q\left( \omega^{1} \right) \right)\left( T + Q\left( \omega^{2} \right) \right)\ldots\left( T + Q\left( \omega^{n} \right) \right)
\end{array}\] \phantomsection\label{permcheck-poly}{}

to be true, in the sense that both sides are the same polynomial in a
single formal variable \(T\). And for that, it is sufficient that a
single random challenge \(T = \lambda\) passes
\hyperref[permcheck-poly]{{[}permcheck-poly{]}}: If the two sides of
\hyperref[permcheck-poly]{{[}permcheck-poly{]}} are not the same
polynomial, then the two sides can have at most \(n - 1\) common values.
So for a randomly chosen \(\lambda\) (chosen from a field with
\(q \approx 2^{256}\) elements), the chances that \(T = \lambda\) passes
\hyperref[permcheck-poly]{{[}permcheck-poly{]}} are extremely small.

We can then get a proof of
\hyperref[permcheck-poly]{{[}permcheck-poly{]}} using the technique of
adding an \emph{accumulator polynomial}. The idea is this: Victor picks
a random challenge \(\lambda \in {\mathbb{F}}_{q}\). Peggy then
interpolates the polynomial
\(F_{P} \in {\mathbb{F}}_{q}\lbrack T\rbrack\) such that
\[\begin{aligned}
F_{P}\left( \omega^{1} \right) & = \lambda + P\left( \omega^{1} \right) \\
F_{P}\left( \omega^{2} \right) & = \left( \lambda + P\left( \omega^{1} \right) \right)\left( \lambda + P\left( \omega^{2} \right) \right) \\
 & \vdots \\
F_{P}\left( \omega^{n} \right) & = \left( \lambda + P\left( \omega^{1} \right) \right)\left( \lambda + P\left( \omega^{2} \right) \right)\cdots\left( \lambda + P\left( \omega^{n} \right) \right).
\end{aligned}\] Then the accumulator
\(F_{Q} \in {\mathbb{F}}_{q}\lbrack T\rbrack\) is defined analogously.

So to prove \hyperref[permcheck-poly]{{[}permcheck-poly{]}}, the
following algorithm works:

strong(Algorithm):\\
Permutation check: Suppose Peggy has committed \(\operatorname{Com}(P)\)
and \(\operatorname{Com}(Q)\).

\begin{enumerate}
\item
  Victor sends a random challenge \(\lambda \in {\mathbb{F}}_{q}\).
\item
  Peggy interpolates polynomials \(F_{P}\lbrack T\rbrack\) and
  \(F_{Q}\lbrack T\rbrack\) such that
  \(F_{P}\left( \omega^{k} \right) = \prod_{i \leq k}\left( \lambda + P\left( \omega^{i} \right) \right)\).
  Define \(F_{Q}\) similarly. Peggy sends \(\operatorname{Com}(F_{P})\)
  and \(\operatorname{Com}(F_{Q})\).
\item
  Peggy uses {[}root-check{]} to prove all of the following statements:

  \begin{itemize}
  \item
    \(F_{P}(X) - \left( \lambda + P(X) \right)\) vanishes at
    \(X = \omega\);
  \item
    \(F_{P}(\omega X) - \left( \lambda + P(\omega X) \right)F_{P}(X)\)
    vanishes at \(X \in \left\{ \omega,\ldots,\omega^{n - 1} \right\}\);
  \item
    The previous two statements also hold with \(F_{P}\) replaced by
    \(F_{Q}\);
  \item
    \(F_{P}(X) - F_{Q}(X)\) vanishes at \(X = 1\).
  \end{itemize}
\end{enumerate}

\subsubsection{Copy check}

Moving on to copy check, let us look at a concrete example where
\(n = 4\). Suppose that our copy constraints were
\[ \redoval{(a_{1})} = 
\redoval{(a_{4})} = 
\redoval{(c_{4})}
\quad\text{and }\quad \bluerect{(b_{2})} = \bluerect{(c_{1})}.\]
(We have colored and circled the variables that will move around for
readability.) So, the copy constraint means we want the following
equality of matrices: \[\begin{pmatrix}
a_{1} & b_{1} & c_{1} \\
a_{2} & b_{2} & c_{2} \\
a_{3} & b_{3} & c_{3} \\
a_{4} & b_{4} & c_{4}
\end{pmatrix} = \begin{pmatrix}
\redoval{(a_{4})} & b_{1} & \bluerect{(b_{2})} \\
a_{2} & \bluerect{(c_{1})} & c_{2} \\
a_{3} & b_{3} & c_{3} \\
\redoval{(c_{4})} & b_{4} & \redoval{(a_{1})}
\end{pmatrix}.\] \phantomsection\label{copy1}{}

Again, our goal is to make this into a \emph{single} equation. There is
a really clever way to do this by tagging each entry with
\(+ \eta^{j}\omega^{k}\mu\) in reading order for \(j = 0,1,2\) and
\(k = 1,\ldots,n\); here \(\eta \in {\mathbb{F}}_{q}\) is any number
such that \(\eta^{2}\) does not happen to be a power of \(\omega\), so
all the tags are distinct. Specifically, if
\hyperref[copy1]{{[}copy1{]}} is true, then for any
\(\mu \in {\mathbb{F}}_{q}\), we also have \[\begin{aligned}
 & \begin{pmatrix}
a_{1} + \omega^{1}\mu & b_{1} + \eta\omega^{1}\mu & c_{1} + \eta^{2}\omega^{1}\mu \\
a_{2} + \omega^{2}\mu & b_{2} + \eta\omega^{2}\mu & c_{2} + \eta^{2}\omega^{2}\mu \\
a_{3} + \omega^{3}\mu & b_{3} + \eta\omega^{3}\mu & c_{3} + \eta^{2}\omega^{3}\mu \\
a_{4} + \omega^{4}\mu & b_{4} + \eta\omega^{4}\mu & c_{4} + \eta^{2}\omega^{4}\mu
\end{pmatrix} \\
 = & \begin{pmatrix}
\redoval{(a_{4})} + \omega^{1}\mu & b_{1} + \eta\omega^{1}\mu & \bluerect{(b_{2})} + \eta^{2}\omega^{1}\mu \\
a_{2} + \omega^{2}\mu & \bluerect{(c_{1})} + \eta\omega^{2}\mu & c_{2} + \eta^{2}\omega^{2}\mu \\
a_{3} + \omega^{3}\mu & b_{3} + \eta\omega^{3}\mu & c_{3} + \eta^{2}\omega^{3}\mu \\
\redoval{(c_{4})} + \omega^{4}\mu & b_{4} + \eta\omega^{4}\mu & \redoval{(a_{1})} + \eta^{2}\omega^{4}\mu
\end{pmatrix}.
\end{aligned}\] \phantomsection\label{copy2}{}

Now how can the prover establish \hyperref[copy2]{{[}copy2{]}}
succinctly? The answer is to run a permutation check on the \(3n\)
entries of \hyperref[copy2]{{[}copy2{]}}! The prover will simply prove
that the twelve matrix entries of the matrix on the left are a
permutation of the twelve matrix entries of the matrix on the right.

The reader should check that this is correct! If the prover starts with
values \(a_{i}\), \(b_{i}\), and \(c_{i}\) that do not satisfy all the
copy constraints, then a randomly selected \(\mu\) is very unlikely to
satisfy this permutation check. The right-hand side will not be a
permutation of the left-hand side, and the check will fail.

To clean things up, shuffle the \(12\) terms on the right-hand side of
\hyperref[copy2]{{[}copy2{]}} so that each variable is in the cell it
started at. We want to prove \[\begin{array}{r}
\begin{pmatrix}
a_{1} + \omega^{1}\mu & b_{1} + \eta\omega^{1}\mu & c_{1} + \eta^{2}\omega^{1}\mu \\
a_{2} + \omega^{2}\mu & b_{2} + \eta\omega^{2}\mu & c_{2} + \eta^{2}\omega^{2}\mu \\
a_{3} + \omega^{3}\mu & b_{3} + \eta\omega^{3}\mu & c_{3} + \eta^{2}\omega^{3}\mu \\
a_{4} + \omega^{4}\mu & b_{4} + \eta\omega^{4}\mu & c_{4} + \eta^{2}\omega^{4}\mu
\end{pmatrix} \\
\text{is a permutation of } \\
\begin{pmatrix}
a_{1} + \redoval{({\eta^{2}\omega^{4}\mu})} & b_{1} + \eta\omega^{1}\mu & c_{1} + \bluerect{({\eta\omega^{2}\mu})} \\
a_{2} + \omega^{2}\mu & b_{2} + \bluerect{({\eta^{2}\omega^{1}\mu})} & c_{2} + \eta^{2}\omega^{2}\mu \\
a_{3} + \omega^{3}\mu & b_{3} + \eta\omega^{3}\mu & c_{3} + \eta^{2}\omega^{3}\mu \\
a_{4} + \redoval{({\omega^{1}\mu})} & b_{4} + \eta\omega^{4}\mu & c_{4} + \redoval{({\omega^{4}\mu})}
\end{pmatrix}.
\end{array}\] \phantomsection\label{copy3}{}

The permutations needed are part of the problem statement, hence
globally known. So in this example, both parties are going to
interpolate cubic polynomials \(\sigma_{a},\sigma_{b},\sigma_{c}\) that
encode the weird coefficients row-by-row: \[\begin{pmatrix}
\sigma_{a}\left( \omega^{1} \right) = \redoval{({\eta^{2}\omega^{4}})} & \sigma_{b}\left( \omega^{1} \right) = \eta\omega^{1} & \sigma_{c}\left( \omega^{1} \right) = \bluerect{({\eta\omega^{2}})} \\
\sigma_{a}\left( \omega^{2} \right) = \omega^{2} & \sigma_{b}\left( \omega^{2} \right) = \bluerect{({\eta^{2}\omega^{1}})} & \sigma_{c}\left( \omega^{2} \right) = \eta^{2}\omega^{2} \\
\sigma_{a}\left( \omega^{3} \right) = \omega^{3} & \sigma_{b}\left( \omega^{3} \right) = \eta\omega^{3} & \sigma_{c}\left( \omega^{3} \right) = \eta^{2}\omega^{3} \\
\sigma_{a}\left( \omega^{4} \right) = \redoval{(\omega^{1})} & \sigma_{b}\left( \omega^{4} \right) = \eta\omega^{4} & \sigma_{c}\left( \omega^{4} \right) = \redoval{(\omega^{4})}.
\end{pmatrix}\] Then the prover can start defining accumulator
polynomials, after re-introducing the random challenge \(\lambda\) from
permutation check. We are going to need six in all, three for each side
of \hyperref[copy3]{{[}copy3{]}}: We call them \(F_{a}\), \(F_{b}\),
\(F_{c}\), \(F'_{a}\), \(F'_{b}\), \(F'_{c}\). The ones on the left-hand
side are interpolated so that \[\begin{aligned}
F_{a}\left( \omega^{k} \right) & = \prod_{i \leq k}\left( a_{i} + \omega^{i}\mu + \lambda \right) \\
F_{b}\left( \omega^{k} \right) & = \prod_{i \leq k}\left( b_{i} + \eta\omega^{i}\mu + \lambda \right) \\
F_{c}\left( \omega^{k} \right) & = \prod_{i \leq k}\left( c_{i} + \eta^{2}\omega^{i}\mu + \lambda \right)
\end{aligned}\] \phantomsection\label{copycheck-left}{}

while the ones on the right have the extra permutation polynomials
\[\begin{aligned}
F'_{a}\left( \omega^{k} \right) & = \prod_{i \leq k}\left( a_{i} + \sigma_{a}\left( \omega^{i} \right)\mu + \lambda \right) \\
F'_{b}\left( \omega^{k} \right) & = \prod_{i \leq k}\left( b_{i} + \sigma_{b}\left( \omega^{i} \right)\mu + \lambda \right) \\
F'_{c}\left( \omega^{k} \right) & = \prod_{i \leq k}\left( c_{i} + \sigma_{c}\left( \omega^{i} \right)\mu + \lambda \right).
\end{aligned}\] \phantomsection\label{copycheck-right}{}

And then we can run essentially the algorithm from before. There are six
initialization conditions \[\begin{aligned}
F_{a}\left( \omega^{1} \right) & = A\left( \omega^{1} \right) + \omega^{1}\mu + \lambda \\
F_{b}\left( \omega^{1} \right) & = B\left( \omega^{1} \right) + \eta\omega^{1}\mu + \lambda \\
F_{c}\left( \omega^{1} \right) & = C\left( \omega^{1} \right) + \eta^{2}\omega^{1}\mu + \lambda \\
F'_{a}\left( \omega^{1} \right) & = A\left( \omega^{1} \right) + \sigma_{a}\left( \omega^{1} \right)\mu + \lambda \\
F'_{b}\left( \omega^{1} \right) & = B\left( \omega^{1} \right) + \sigma_{b}\left( \omega^{1} \right)\mu + \lambda \\
F'_{c}\left( \omega^{1} \right) & = C\left( \omega^{1} \right) + \sigma_{c}\left( \omega^{1} \right)\mu + \lambda.
\end{aligned}\] \phantomsection\label{copycheck-init}{}

and six accumulation conditions \[\begin{aligned}
F_{a}(\omega X) & = F_{a}(X) \cdot \left( A(\omega X) + X\mu + \lambda \right) \\
F_{b}(\omega X) & = F_{b}(X) \cdot \left( B(\omega X) + \eta X\mu + \lambda \right) \\
F_{c}(\omega X) & = F_{c}(X) \cdot \left( C(\omega X) + \eta^{2}X\mu + \lambda \right) \\
F'_{a}(\omega X) & = F'_{a}(X) \cdot \left( A(\omega X) + \sigma_{a}(X)\mu + \lambda \right) \\
F'_{b}(\omega X) & = F'_{b}(X) \cdot \left( B(\omega X) + \sigma_{b}(X)\mu + \lambda \right) \\
F'_{c}(\omega X) & = F'_{c}(X) \cdot \left( C(\omega X) + \sigma_{c}(X)\mu + \lambda \right)
\end{aligned}\] \phantomsection\label{copycheck-accum}{}

before the final product condition
\[F_{a}(1)F_{b}(1)F_{c}(1) = F'_{a}(1)F'_{b}(1)F'_{c}(1).\]
\phantomsection\label{copycheck-final}{}

To summarize, the copy check goes as follows:

strong(Algorithm):\\
Copy check:

\begin{enumerate}
\setcounter{enumi}{-1}
\item
  Peggy has already sent the three commitments
  \(\operatorname{Com}(A),\operatorname{Com}(B),\operatorname{Com}(C)\)
  to Victor; these commitments bind her to the values of all the
  variables \(a_{i}\), \(b_{i}\), and \(c_{i}\).
\item
  Both parties compute the degree \(n - 1\) polynomials
  \(\sigma_{a},\sigma_{b},\sigma_{c} \in {\mathbb{F}}_{q}\lbrack X\rbrack\)
  described above, based on the copy constraints in the problem
  statement.
\item
  Victor chooses random challenges \(\mu,\lambda \in {\mathbb{F}}_{q}\)
  and sends them to Peggy.
\item
  Peggy interpolates the six accumulator polynomials \(F_{a}\), \ldots,
  \(F'_{c}\) defined in \hyperref[copycheck-left]{{[}copycheck-left{]}}
  and \hyperref[copycheck-right]{{[}copycheck-right{]}}.
\item
  Peggy uses {[}root-check{]} to prove
  \hyperref[copycheck-init]{{[}copycheck-init{]}} holds.
\item
  Peggy uses {[}root-check{]} to prove
  \hyperref[copycheck-accum]{{[}copycheck-accum{]}} holds for
  \(X \in \left\{ \omega,\omega^{2},\ldots,\omega^{n - 1} \right\}\).
\item
  Peggy uses {[}root-check{]} to prove
  \hyperref[copycheck-final]{{[}copycheck-final{]}} holds.
\end{enumerate}
